% !TEX program = xelatex
%==============================================================================
%  宋睿 · 简历 —— 投递「中铝乾星（成都）科技有限责任公司 · 工艺技术主管」
%
%  编译：TeXstudio 按 F5　／　命令行 xelatex resume.tex（跑两遍）
%
%  文件结构分两层：
%    ┌ 排版层：从下面到 "CONTENT 内容层" 标记之间 —— 字体、颜色、间距、宏
%    └ 内容层："CONTENT 内容层" 标记之后 —— 本次投递的实际内容
%==============================================================================
\documentclass[11pt,a4paper]{ctexart}

%------------------------------------------------------------------------------
% 0. 让 PDF 里的中文可被复制、被招聘系统（ATS）正确解析。【勿删】
%------------------------------------------------------------------------------
\XeTeXgenerateactualtext=1

%------------------------------------------------------------------------------
% 1. 字体：正文宋体，加粗自动切黑体
%------------------------------------------------------------------------------
\usepackage{fontspec}
\setCJKmainfont{Songti SC}[BoldFont=Heiti SC]
\setCJKsansfont{Heiti SC}
\newCJKfontfamily\hei{Heiti SC}      % \hei  切黑体（章节标题用）
\newCJKfontfamily\song{Songti SC}    % \song 切宋体
\newCJKfontfamily\kai{Kaiti SC}      % \kai  切楷体
\newCJKfontfamily\songb{Songti SC Bold}  % \songb 宋体加粗（条目主标题用）

%------------------------------------------------------------------------------
% 2. 页面与基础排版
%------------------------------------------------------------------------------
\usepackage[top=1.15cm,bottom=0.95cm,left=1.45cm,right=1.45cm]{geometry}
\usepackage{xcolor}
\usepackage{graphicx}
\usepackage{tabularx}
\usepackage{array}
\usepackage{enumitem}
\usepackage{etoolbox}
\usepackage[hidelinks]{hyperref}

\pagestyle{empty}
\setlength{\parindent}{0pt}
\linespread{1.0}
\setlength{\parskip}{0pt}

\definecolor{accent}{HTML}{1F4E79}   % 主题色（深蓝）
\definecolor{rulegray}{HTML}{2B2B2B}

%------------------------------------------------------------------------------
% 3. 排版宏
%------------------------------------------------------------------------------

% 章节标题：\ressection{教育背景}
%   上方留白大于下方，使「标题 + 横线」读作一个整体
\newcommand{\ressection}[1]{%
  \vspace{0.52em}%
  \noindent{\hei\large\bfseries\color{accent}#1}\par
  \vspace{-0.55em}%
  \noindent{\color{accent}\rule{\linewidth}{1pt}}\par
  \vspace{0.10em}%
}

% 条目头：\entry{机构或项目名}{时间}{副标题}
%   整行统一宋体：主标题宋体加粗，副标题宋体常规且略小，时间右对齐
%   不需要副标题就写 {}
\newcommand{\entry}[3]{%
  \vspace{0.14em}%
  \noindent\begin{tabularx}{\linewidth}{@{}X@{}r@{}}
    {\songb #1}\ifblank{#3}{}{\hspace{0.7em}{\small #3}} & {\small #2} \\
  \end{tabularx}%
}

% 要点列表：topsep 与 itemsep 取同一值，保证「条目→首条」与「条→条」间距一致
\newenvironment{bullets}{%
  \begin{itemize}[leftmargin=1.0em,itemsep=0.10em,topsep=0.10em,parsep=0pt,%
                  partopsep=0pt,label=\textbullet]%
}{\end{itemize}}

% 技能行：\skillrow{类别}{内容}
\newcommand{\skillrow}[2]{%
  \noindent\textbf{#1}：#2\par\vspace{0.04em}%
}

% 页眉：\resumeheader{姓名}{联系方式}{照片}{求职意向}　照片不存在会自动画占位框
\newcommand{\resumeheader}[4]{%
  \noindent
  \begin{minipage}[c]{\dimexpr\linewidth-2.75cm\relax}
    {\hei\fontsize{22}{26}\selectfont #1}\par
    \ifblank{#4}{}{\vspace{0.28em}{\small\bfseries\color{accent}求职意向：#4}\par}%
    \vspace{0.34em}
    {\small #2}\par
  \end{minipage}%
  \hfill
  \begin{minipage}[c]{2.5cm}
    \centering
    \IfFileExists{#3}%
      {\includegraphics[width=2.2cm,height=2.9cm,keepaspectratio]{#3}}%
      {\fbox{\parbox[c][2.9cm][c]{2.2cm}{\centering\scriptsize 照片\\[0.2em]放\\[0.2em]\texttt{\tiny #3}}}}%
  \end{minipage}\par
  \vspace{0.30em}%
}

% 压缩开关：内容溢出时在内容层开头加 \tighten
\newcommand{\tighten}{%
  \linespread{1.0}\selectfont
  \setlist[itemize]{itemsep=0.0em,topsep=0.02em,parsep=0pt}%
  \setlength{\parskip}{0pt}%
}

%==============================================================================
%  CONTENT 内容层
%==============================================================================
% 内容溢到第二页时去掉下面这行的注释（会把行距和列表间距全局压紧）
% \tighten

\begin{document}

\resumeheader{宋　睿}%
  {13298121974~|~songrui0006@163.com~|~中共党员~|~女~|~2000.06}%
  {photo.jpg}%
  {工艺技术（材料 / 化工方向）}

%------------------------------------------------------------------ 教育背景
\ressection{教育背景}

\entry{中国海洋大学}{2024.09 -- 2027.06}{化学工程与技术（硕士）~|~GPA 3.58}
\begin{bullets}
  \item \textbf{主修课程}：分子模拟计算（90）、高等仪器分析（93）、数值分析（96）
\end{bullets}

\entry{河南工业大学}{2018.09 -- 2022.07}{化学工程与工艺（本科）~|~GPA 3.76（专业前 10\%）}
\begin{bullets}
  \item \textbf{主修课程}：物理化学（95）、化工原理（91）、高等数学（96）、实验设计与数据处理（89）
\end{bullets}

%------------------------------------------------------------------ 科研经历
\ressection{科研经历}

\entry{电催化硝酸根还原合成氨（eNO$_3$RR）}{2025.12 -- 至今}{中国科学院青岛生物能源与过程研究所}
\begin{bullets}
  \item 设计铜基催化剂，通过缺陷与界面工程调控电子结构；系统筛选不同形貌 Cu 并引入 Co(OH)$_2$ 改性，使氨产率提升 7 倍、法拉第效率接近 100\%
  \item 建立质谱原位监测方法追踪反应路径，结合 DFT 计算定位活性位点，形成「合成—测试—机理」完整闭环
  \item 成果：论文撰写中
\end{bullets}

\entry{电催化氮气氧化合成硝酸根（eNOR）}{2024.08 -- 2025.11}{中国海洋大学}
\begin{bullets}
  \item 制备 $\delta$-MnO$_2$/MoS$_2$ 异质结催化剂，调控组分比例筛选最优配比；引入过一硫酸氢钾（PMS）高级氧化工艺提升性能
  \item 产率达 198.5 $\mu$g·h$^{-1}$·mg$^{-1}$，法拉第效率 19.31\%；结合原位拉曼与 DFT 计算揭示反应机理
  \item 成果：论文在投
\end{bullets}

\entry{电化学水凝胶传感器（H$_2$O$_2$ 检测）}{2024.09 -- 2025.01}{中国海洋大学}
\begin{bullets}
  \item 自主合成导电水凝胶材料并修饰电极，运用循环伏安法系统优化检测性能
\end{bullets}

%------------------------------------------------------------------ 项目经历
\ressection{项目经历}

\entry{海藻酸钠水凝胶的制备及吸附性能研究}{2019.12 -- 2022.07}{项目负责人~|~河南工业大学 党丹老师课题组}
\begin{bullets}
  \item 以离子交联法制备水凝胶，系统研究交联条件与吸附环境对性能的影响；优化条件下对刚果红去除率达 95\%，理论最大吸附量 180 mg/g
  \item 结合吸附等温与动力学模型拟合及 ESEM、BET、FTIR 表征，阐明吸附机理
  \item 统筹 10 人以上团队并培养 8 名低年级组成员；成果发表于 SCI 期刊并获授权专利
\end{bullets}

\entry{年产 10 万吨异丙醇工艺设计}{2021.02 -- 2021.07}{第十五届全国大学生化工设计大赛~|~小组成员}
\begin{bullets}
  \item 经文献调研选定丙酮催化加氢为核心工艺，结合市场、政策与供应链确定经济生产规模
  \item 独立完成可行性研究报告、经济评价、「三废一噪」处理及 HSE 分析全套技术文档，形成完整工业化设计方案；获华中赛区三等奖
\end{bullets}

%------------------------------------------------------------------ 学生工作
\ressection{学生工作}

\skillrow{研究生兼职辅导员（2025.03 -- 2025.08）}%
  {中国海洋大学化学化工学院。参与党建工作全流程，承担会务协助、信息整理、文档校对、考核评奖与档案管理}
\skillrow{党支部组织委员兼纪检委员（2024.08 -- 至今）}%
  {2024 级硕士研究生第三支部。组织主题党日活动，负责党员信息库动态更新与党费收缴公示}

%------------------------------------------------------------------ 专业技能
\ressection{专业技能}

\skillrow{材料合成与制备}{异质结催化剂合成、缺陷与界面工程调控、离子交联、自由基聚合、生物炭负载}
\skillrow{表征与测试}{循环伏安等电化学测试、原位拉曼、原位质谱、ESEM、BET、FTIR}
\skillrow{计算与数据分析}{DFT 计算、分子模拟、吸附等温与动力学模型拟合、实验设计与数据处理}
\skillrow{文档与办公}{熟练使用常用办公软件；独立撰写过可行性研究报告、经济评价、HSE 分析等工程技术文档}

%------------------------------------------------------------------ 荣誉与成果
\ressection{荣誉与成果}

\skillrow{学术成果}{SCI 论文 1 篇（Adv.\ Mater.\ Sci.\ Eng., 2023）、授权专利 1 项；硕士期间论文在投、在写各 1 篇}
\skillrow{竞赛与荣誉}{「互联网+」大赛河南赛区一等奖、全国化工设计大赛华中赛区三等奖；优秀研究生、优秀兼职辅导员；校优秀学生奖学金 2 次、优秀共青团员 2 次}

\end{document}
